Artificial intelligence now occupies a central place in corporate disclosure, yet the informational content of that language remains difficult to evaluate ex ante. AI-related investments have real implications for product innovation, operating efficiency, and firm growth (Babina et al., 2024; Cao et al., 2024), and investors appear to respond to AI narratives in corporate filings rather than treating them as generic boilerplate (Basnet et al., 2025). At the same time, the recent expansion of AI-related language in annual reports has outpaced the ability of outside investors to verify which disclosures reflect implemented technological capability and which primarily reflect narrative positioning. In our sample, more than 40\% of AI-talking firm-years in both 2024 and 2025 are classified as PatentMismatch firm-years, meaning that the filing language appears low in credibility while contemporaneous AI patenting remains weak. That pattern is difficult to reconcile with a view in which the recent surge in AI disclosure is, on average, a transparent reflection of the underlying implementation.

The fundamental economic issue is not unique to AI. In disclosure settings where implementation is difficult to verify in real time, communication can be informative without being tightly constrained by realized action. That logic is familiar from both disclosure theory and the empirical washing literature. Signaling models emphasize that communication is informative only when lower-quality senders cannot cheaply mimic higher-quality ones (Spence, 1973; Seidmann \& Winter, 1997), while empirical work in adjacent domains shows that public claims can move ahead of later observable conduct when the topic is salient and externally valued (Delmas \& Burbano, 2011; Kim \& Lyon, 2015; Lyon \& Montgomery, 2015; Baker et al., 2024). Clearly, AI-related disclosures belong to that class of problems. The underlying activity is technologically consequential, highly visible, and often difficult for investors to verify contemporaneously, making the credibility of the narrative itself an empirical question rather than a rhetorical one.

Recent work has begun to study AI communication more directly, and the closest papers provide the immediate empirical starting point for the present design. Basnet et al.~(2025) show that investors respond differently to actionable and speculative AI language in 10-K business descriptions, while Cao et al.~(2024) show that AI disclosure contains forward-looking information about firm outcomes. Bandyopadhyay et al.~(2023) connect broader AI narratives to subsequent stock mispricing, especially among firms outside the largest size group, and more recent working-paper evidence studies the divergence between AI talk and AI walk more explicitly and traces that divergence to market pricing, institutional trading, and managerial incentives (Li, 2026; Donelson et al., 2025; Barrios et al., 2025). The open question is how far a filing-based credibility measure can be scaled and audited while still preserving the distinction between more credible and less credible AI disclosure, and whether that same construct can be taken jointly to market-pricing outcomes and later technological realization.

That unresolved problem is as much about measurement as it is about economic interpretation. The textual-analysis literature has shown for some time that corporate language can contain information about later fundamentals, returns, and operating outcomes, but it has also shown how sensitive these inferences are to the way language is measured. Forward-looking statements in filings predict subsequent fundamentals (Li, 2010), disclosure language contains information beyond reported numbers (Davis et al., 2012), and textual measures in finance require domain-specific construction rather than imported dictionaries or generic sentiment tools (Loughran \& McDonald, 2011, 2016). Within innovation settings, richer text measures capture meaningful variation in firm-level technological activity and forecast subsequent patenting and performance (Bellstam et al., 2021), while public disclosure improves how markets interpret innovation-related assets (Kim \& Valentine, 2023). The question here is therefore not simply whether firms mention AI more often. It is whether the composition of that language isolates a subset of disclosures that are systematically less aligned with contemporaneous implementation and later technological realization.

Against that backdrop, we study annual 10-K filings of U.S. public firms from 2016 to 2025 and classify AI-related sentences as actionable, speculative, or irrelevant. The classification is produced by a scalable audited pipeline rather than by manual coding alone, and the resulting sentence-level judgments are aggregated into firm-year measures of disclosure composition, including the AS ratio and PatentMismatch. The latter combines low-credibility disclosure composition with weak contemporaneous AI patenting, thereby isolating a subset of AI-talking firm-years in which the narrative is difficult to reconcile with contemporaneous technological output. That same disclosure layer is then linked both to subsequent AI patent outcomes and to market measures around and after the filing date. The design is deliberately joint: if the construct captures something real, it should be informative about later innovation, and if markets do not fully distinguish credible from non-credible disclosure at the filing date, that same construct should matter for the subsequent return path.

The empirical evidence follows a clear sequence. It begins with the sharp post-2022 rise in PatentMismatch and characterizes 41.5\% of AI-talking firm-years in 2024 and 42.1\% in 2025. However, filing-date abnormal returns provide little evidence of immediate discipline, whereas the sharper pricing pattern appears after disclosure: a long--short strategy that buys firms with more credible AI disclosure and shorts mismatch firms earns its strongest return over the subsequent quarter, with an equal-weight alpha of 0.93\% per month, or roughly 11.2\% annualized. That effect is much weaker in value-weighted specifications, suggesting that the slower correction is concentrated outside the largest firms. In that respect, the return evidence is closer to delayed incorporation than to immediate repricing and is consistent with filing-based settings in which disclosure-related information is not fully impounded at the announcement date (Cohen et al., 2020). Furthermore, the disclosure construct remains tied to later technological realization: actionable AI language aligns more closely with subsequent AI patenting than speculative language, and the AS-ratio $\times$ PatentMismatch specification predicts weaker subsequent AI patent output. The post-ChatGPT tests point in the same direction. What makes the pattern economically important is not only that low-credibility AI disclosure became widespread, but that the market appears to underreact to it even as it predicts weaker subsequent technological realization.

Basnet et al.~(2025) provide the closest filing-based starting point by showing that investors distinguish between actionable and speculative AI language in 10-K business descriptions, while Bandyopadhyay et al.~(2023) show that broader AI narratives are associated with subsequent stock mispricing, especially outside the largest firms. Li (2026) moves even closer to the present setting by studying the divergence between AI talk and AI walk and its market consequences. Building on those studies, the present paper scales the actionable--speculative distinction into a reusable audited framework for mandatory annual filings, introduces PatentMismatch as a filing-based construct that combines disclosure composition with contemporaneous patent weakness, and uses that construct to connect low-credibility AI disclosure both to slower market correction and to weaker subsequent AI patent realization. The contribution therefore lies in identifying which subset of filing-based AI narratives appears least credible and in showing that this subset is economically important in both prices and later technological realization.

The remainder of the paper proceeds as follows. Section 2 reviews the related literature. Section 3 describes the data, the measurement pipeline, and the empirical design. Section 4 presents the main results. Section 5 concludes.